
\subsubsection{Principal Findings}

\begin{enumerate}
\item \textbf{Sample efficiency gap:} NFN achieves $R^{2}=0.95$ at N=1,000 where MLP achieves $R^{2}=0.35$, a 60 percentage point difference.
\end{enumerate}

\begin{enumerate}
\item \textbf{NFN invariance is near-perfect:} Deviation below 1.$19\times 10^{-7}$ across permutations confirms mathematical guarantee.
\end{enumerate}

\begin{enumerate}
\item \textbf{MLP dissociation:} $R^{2}=0.986$ coexists with invariance=0.63 at N=40,000. Task performance and mechanism quality are separable.
\end{enumerate}

\begin{enumerate}
\item \textbf{Invariance cannot be learned:} MLP invariance at N=40K (0.63) is not meaningfully higher than at N=1K (0.64), despite $R^2$ improving from 0.004 to 0.986.
\end{enumerate}

\subsubsection{Interpretation}

MLPs trained on large-scale weight-space data learn to predict accuracy by memorizing position-sensitive statistics---which weight indices correlate with accuracy in the specific dataset. This strategy achieves high $R^2$ when training and test distributions align but does not generalize the underlying permutation symmetry.

NFN's architectural equivariance extracts accuracy-predictive features from a canonical view of each model's weights, enabling high performance with limited data. The inductive bias is structural, not compensable by data quantity.

\subsubsection{H-M4 Interpretation}

H-M4 showed MLP invariance of 0.64 at N=1K, which exceeded the 0.5 threshold. However, the corresponding $R^2$ was 0.004, indicating the model had not learned the task. The apparent ''high invariance'' was an artifact of near-constant predictions. This highlights that invariance metrics should be conditioned on minimum task performance.

\subsubsection{Limitations}

\begin{enumerate}
\item \textbf{Synthetic data in H-E1:} The initial proof-of-concept (H-E1) used synthetic MLP weights rather than real Model Zoo data. Subsequent hypotheses (H-M4, H-M5) used real Model Zoo data, confirming the overall pattern.
\end{enumerate}

\begin{enumerate}
\item \textbf{Single architecture family:} All experiments used CIFAR-10 CNNs. Generalization to transformers or other architectures is untested.
\end{enumerate}

\begin{enumerate}
\item \textbf{NFN-Scrambled baseline not implemented:} Prediction P2 (comparing NFN with correct versus incorrect permutation group) was not tested.
\end{enumerate}

\begin{enumerate}
\item \textbf{Limited seeds:} Some experiments used 1--3 seeds rather than the planned 10. Effect sizes were large enough that conclusions are unlikely to change with additional seeds.
\end{enumerate}

\subsubsection{Broader Implications}

These results suggest that certain symmetries require explicit architectural encoding; data quantity cannot substitute for structural inductive bias. This principle may extend beyond weight-space learning to other domains where input symmetries are present but not architecturally exploited.

